\section{Introduction}

\subsection{The reproducibility crisis}

Modern science is facing what has been called the ``reproducibility crisis,'' where scientists are unable to reproduce each others' experiments \cite{natureRepro,manifesto4ReproducibleScience}.
Lack of reproducibility threatens two of the four Mertonian norms \cite{sociologyOfScience} of science: communalism, as one may no longer build off of anothers' development, and disinterested skepticism, as one may no longer scrutinize anothers' finding.
With these norms in question, science fails to self-correct and, if persists long enough, may fail to command trustworthiness.
There have been many high-profile erroneous conclusions that were uncovered by failure to reproduce the results \cite{herndonCritiqueReinhartRogoff,scientistNightmareSoftware,fMRIClusterFailure,neupaneRetraction}.

Scientists polled gave several reasons for irreproducibility \cite{natureRepro} that generally fall into three overlapping categories:
\begin{itemize}
  \item \textbf{misaligned incentives}: e.g., pressure to publish, positive results bias, less prestige for reproducing
  \item \textbf{methodological issues}: e.g., poor experimental design, spurious correlations, low statistical power, methods or data not available, computation
  \item \textbf{fraud}
\end{itemize}

One set of methodological issues are the difficulties encountered when reproducing a computation.
Unlike arbitrary physical systems, conventional computational systems are digital, so it is feasible to perfectly describe and perfectly recreate their states.
Perfect description and recreation \textit{should} make reproducing computational experiments much simpler, but practices misses the mark.
In practice, irreproducibility is often due to authors not sharing relevant code, data, or dependency specifications \cite{reproducibleSignalProc,whyWorkflowsBreak,collbergTechRep}.

\todo[Why is code missing?]

As more scientists use computational methods, software irreproducibility will affect more fields.
Of university scientists polled, 90\% use research software in their research and 70\% say that their research would not be practical without it \cite{researchSoftwareSurvey}.
At the same time, 56\% develop their own software in the process of their research \cite{researchSoftwareSurvey}.
Any purported solution to software irreproducibility must be deployable by many different research labs.


Research software should be reproducible

Reasons for non-reproducibility

\subsection{Scope of this work}
